\begin{abstract}
The gLOWCOST detector uses plastic scintillator tiles, wavelength-shifting fibres, and silicon photomultipliers (SiPMs) to detect charged particles. Stable operation requires an SiPM bias that is high enough to detect scintillation light but not so high that dark counts, correlated avalanches, or readout saturation dominate. This work develops and tests a waveform-based method for measuring the breakdown voltage of Hamamatsu S13360-2050VE SiPMs through the existing gLOWCOST analogue readout. Dark-pulse waveforms were recorded with a PicoScope, corrected event by event for baseline offset, and integrated in a peak-aligned time gate. The spacing between neighboring photoelectron populations was used as a relative gain measurement and extrapolated linearly to zero spacing.

The study also establishes an EPIC-tile coincidence method for measuring the full response of GSU gLOWCOST tiles. This zero-event measurement is intended to test whether channels operated at equal overvoltage have equal practical detection efficiency. The current scintillator data are treated as method development rather than a completed photon-detection-efficiency measurement. Once the method has been validated, a pulsed-LED study can provide a true SiPM zero-event calibration.
\end{abstract}
